\subsection{Estructura Interna (Patrón AAA)}

\subsubsection{Regla: separar Arrange, Act y Assert con líneas en blanco}

\textbf{Regla:} Tres bloques consecutivos --- preparación, ejecución, verificación --- separados solo por una línea en blanco, sin comentarios que los etiqueten (Osherove, 2013).

\textbf{Descripción:} \texttt{ExpectEqual} representa la llamada de aserción del bloque de verificación.

\textbf{Código correcto:}
\begin{lstlisting}[style=csharpstyle]
public void Roll_FaceValueSix_GrantsExtraTurn()
{
    DiceRoller diceRoller = new(new FixedRandomProvider(6));

    int faceValue = diceRoller.Roll();

    ExpectEqual(6, faceValue);
}
\end{lstlisting}

\textbf{Código incorrecto:}
\begin{lstlisting}[style=csharpstyle]
public void Roll_FaceValueSix_GrantsExtraTurn()
{
    DiceRoller diceRoller = new(new FixedRandomProvider(6));
    int faceValue = diceRoller.Roll();
    ExpectEqual(6, faceValue);
}
\end{lstlisting}

\subsubsection{Regla: una única aserción por prueba}

\textbf{Regla:} El bloque de verificación contiene una única llamada de aserción, como \texttt{ExpectEqual} (Osherove, 2013). Para comparar varios atributos de un mismo objeto, se prefiere una sola aserción sobre el propio \texttt{Equals} de la clase (Sección 2.9) en lugar de una aserción por atributo. Un caso puntual que en verdad necesite más de una aserción se admite solo con un comentario, inmediatamente antes de las aserciones, que explique el motivo de esa decisión.

\textbf{Código correcto:}
\begin{lstlisting}[style=csharpstyle]
public void MoveTo_ValidDestination_UpdatesCurrentTile()
{
    Player player = new(startingTile: new BoardTile(0, 0));

    player.MoveTo(new BoardTile(2, 3));

    ExpectEqual(new BoardTile(2, 3), player.CurrentTile);
}
\end{lstlisting}

\textbf{Código incorrecto:}
\begin{lstlisting}[style=csharpstyle]
public void MoveTo_ValidDestination_UpdatesCurrentTile()
{
    Player player = new(startingTile: new BoardTile(0, 0));

    player.MoveTo(new BoardTile(2, 3));

    ExpectEqual(2, player.CurrentTile.Row);
    ExpectEqual(3, player.CurrentTile.Column);
}
\end{lstlisting}
